\begin{helper}
Let $k>0$, let $\mathcal F\in H(k,k,D)$, and let $v$ be a vertex with
$\deg_{\mathcal F}(v)<D$.  If
$\operatorname{FreshEdge}_k(\mathcal F,v)$ is the hypergraph obtained by
adjoining one fresh $k$-edge through $v$ as in
\noderef{KUniformFreshEdgeHypergraph}, then
$\operatorname{FreshEdge}_k(\mathcal F,v)$ is also a member of $H(k,k,D)$.
\end{helper}
\begin{proof}
By \noderef{HypergraphClass}, membership in $H(k,k,D)$ means
$k$-uniformity, $k$-simplicity, and maximum degree at most $D$.
The old edges of the fresh-edge hypergraph are unchanged except that their
vertices are viewed inside the enlarged vertex set, while the new edge is
the vertex $v$ together with exactly $k-1$ fresh vertices by
\noderef{KUniformFreshEdgeHypergraph}.  Hence every old edge still has
cardinality $k$ by \noderef{UniformHypergraph}, and the new edge has
$1+(k-1)=k$ vertices because $k>0$.

For simplicity, two old edges have the same intersection size they had in
$\mathcal F$, so their intersection has size at most $k$ by the
\noderef{TSimpleHypergraph} part of the hypothesis.  The new edge meets an
old edge only possibly at the old vertex $v$, since all other vertices of
the new edge are fresh; consequently such an intersection has size at most
$1\le k$.  There is only one new edge label, so there is no distinct
new--new pair to check.

Finally consider degrees, using \noderef{HypergraphDegree}.  A fresh vertex
lies in exactly the new edge, so its degree is $1$, and this is at most $D$
because $\deg_{\mathcal F}(v)<D$ implies $D>0$.  An old vertex different
from $v$ lies in exactly the same old edge labels as before and not in the
new edge, so its degree remains at most $D$ by \noderef{MaxDegreeAtMost}.
The vertex $v$ gains exactly one incident edge, and its old degree was
strictly less than $D$, so its new degree is at most $D$.  Thus all three
parts of \noderef{HypergraphClass} hold for the fresh-edge hypergraph.
\end{proof}
